\chapter{Optimisation Framework}
\label{chap:optimization}

This chapter describes the Optuna-TPE-based optimisation framework
(optimizer/optimize.py) used to search for high-quality inventory
and transport parameters.  Three modes are supported: inventory-only,
transport-only, and joint.  The joint mode is the central research
contribution of this study.

%% -----------------------------------------------------------------------
\section{Need for Optuna TPE}
\label{sec:opt:tpe}

The optimisation problem of selecting base-stock levels and dispatch
thresholds for a multi-echelon network is a \emph{black-box, mixed-integer,
constrained} problem: the objective function (total simulated cost subject
to a fill-rate constraint) has no closed-form gradient, the decision
variables include both integer-valued stock levels and continuous-valued
dispatch thresholds, and the feasible region is defined by a fill-rate
inequality that cannot be expressed analytically.

Classical derivative-based methods are inapplicable.  Genetic algorithms
and simulated annealing are gradient-free but require many evaluations to
converge.  Bayesian optimisation with a Gaussian-process surrogate is
sample-efficient but scales poorly beyond $\sim\!20$ dimensions due to
the cubic cost of kernel inversion.

The \textbf{Tree-structured Parzen Estimator (TPE)} algorithm addresses these limitations:
\begin{itemize}
  \item It models the distribution of good and bad parameters separately
        using kernel density estimators, requiring no global surrogate.
  \item It scales approximately linearly in the number of dimensions,
        making it practical for 25--30 dimensional joint spaces.
  \item It handles mixed-integer variables natively — integer and
        continuous parameters are treated uniformly.
  \item It is implemented in the open-source \textsc{Optuna} library
        with support for parallelism and pruning.
\end{itemize}

%% -----------------------------------------------------------------------
\section{Search-Space Encoding}
\label{sec:opt:space}

\subsection{Inventory: base-stock levels}
\label{sec:opt:space-inventory}

For the inventory mode, the optimiser searches over the base-stock level
$S_{n,s}$ for every non-supplier (node, SKU) pair.  In the 1N3 topology
with $N=2$ warehouses (W1, W2) and 3 retailers (R1, R2, R3) — five
non-supplier nodes in total — and 5 SKUs, this gives $5 \times 5 = 25$
integer parameters.

Each parameter $S_{n,s}$ is bounded relative to the analytical base-stock
level $S^*_{n,s}$ (Equation~\ref{eq:newsvendor}):

\begin{equation}
  S_{n,s} \in \Bigl[\max\!\bigl(10,\; \lfloor 0.3 \cdot S^*_{n,s} \rfloor\bigr),\;
                    \lceil 3.0 \cdot S^*_{n,s} \rceil\Bigr] \cap \mathbb{Z}
  \label{eq:inv-bounds}
\end{equation}

The lower bound of 10 units prevents degenerate stock-out configurations.
The upper bound of $3 \times S^*$ is wide enough to encompass very
conservative safety-stock strategies.

%% -----------------------------------------------------------------------
\subsection{Transport: dispatch thresholds}
\label{sec:opt:space-transport}

For the transport mode, the optimiser searches over the
min\_dispatch\_utilization threshold $\delta_e$ for every edge $e$
in the network.  The 1N3 topology has 5 directed edges
(Supplier$\to$W1, Supplier$\to$W2, W1$\to$R1, W1$\to$R2, W2$\to$R3),
giving 5 continuous parameters:

\begin{equation}
  \delta_e \in [0.0,\; 0.9]
  \label{eq:disp-bounds}
\end{equation}

A value of 0 means dispatch on any positive load; a value of 0.9 means
the planner waits until 90\% of the vehicle's capacity is committed.
The upper bound of 0.9 (rather than 1.0) prevents the edge from
deadlocking when demand is low.

%% -----------------------------------------------------------------------
\section{Joint Search}
\label{sec:opt:joint}

In joint mode, a single Optuna study is created over the \emph{union} of
the inventory and transport search spaces — 30 parameters for the 1N3
($N=2$) topology (25 integer stock levels + 5 continuous dispatch
thresholds).  A single trial proposes all 30 values simultaneously,
evaluates the full simulation, and returns the penalised cost.

This differs from an alternating or decomposed optimisation approach:
both sets of parameters are updated together in each trial, allowing
the TPE model to learn cross-parameter interactions — specifically, that
raising $S_{n,s}$ (inventory) enables raising $\delta_e$ (consolidation)
without increasing stockout risk.

The joint study runs $n_\text{trials}$ sequential trial evaluations by default.

%% -----------------------------------------------------------------------
\section{Constrained Optimisation --- Minimum Fill Rate}
\label{sec:opt:constraint}

Most supply chains operate under a minimum service-level agreement.  The
optimiser enforces a minimum fill rate $\alpha_{\min}$ via a
\emph{soft penalty}:

\begin{equation}
  \text{objective} = C_{\text{total}} + \max\!\bigl(0,\;
    \alpha_{\min} - \hat{\alpha}\bigr) \times 100 \times 10^6
  \label{eq:penalty}
\end{equation}

where $\hat{\alpha}$ is the fill rate achieved by the current trial and
the $10^6$ multiplier per percentage-point shortfall makes a 1\%
fill-rate violation cost \$1 million — far exceeding any realistic cost
saving.  This ensures that feasible trials (those with $\hat{\alpha} \geq
\alpha_{\min}$) always dominate infeasible ones in the Pareto ranking.

After all trials are complete, the best \emph{feasible} trial is selected
as the result.  If no trial satisfies the fill-rate constraint, the
framework falls back to the best-overall trial and warns the user.

\subsection{Aggregate vs per-SKU constraint}
\label{sec:opt:per-sku}

Equation~\ref{eq:penalty} enforces the constraint on the
\emph{aggregate} fill rate, allowing high-volume SKUs to over-serve
and compensate for under-served low-volume SKUs.  In service-level
agreements that guarantee fill rate \emph{per item}, the constraint
should bind for every SKU individually.  The optimiser supports this
via a per\_sku\_constraint=True flag, which replaces the scalar
penalty with a per-SKU sum:

\begin{equation}
  \text{penalty}_{\text{per-SKU}} \;=\; \sum_{s \in \text{SKUs}}
        \max\!\bigl(0,\; \alpha_{\min} - \hat{\alpha}_s\bigr) \times 100 \times 10^6
  \label{eq:penalty-sku}
\end{equation}

A trial is feasible only when every SKU's fill rate meets the target.
Experiment~E3\_per\_sku (Section~\ref{sec:results:e3_per_sku}) re-runs the
joint optimisation under this stricter mode and reports the cost premium.

%% -----------------------------------------------------------------------
\section{Pareto Exploration}
\label{sec:opt:pareto}

Experiment E4 maps the cost--service Pareto frontier by two complementary methods: NSGA-II (primary) and a constraint sweep (sanity check).

\subsection{Primary: NSGA-II two-objective study}

The Optuna NSGAIISampler is run on the same
30-dimensional joint search space, but with two minimisation objectives:

\begin{equation}
  \min \bigl(C_{\text{total}},\; -\hat{\alpha}\bigr)
  \label{eq:nsga2-objective}
\end{equation}

After 200 trials the non-dominated trials in study.best\_trials
form the Pareto front.  This produces a smooth, monotone cost--service
trade-off curve in a single run, without needing to choose fill-rate
target points up-front.  The NSGA-II frontier is the primary result reported in Chapter~\ref{chap:results}.

\subsection{Sanity check: constraint sweep}

The original constraint-sweep procedure is retained as a sanity-check
overlay.  The joint optimiser is re-run at
eight integer-step targets $\alpha_{\min} \in [92\%,\,99\%]$,
each independent and seeded differently.  At lower targets the
constraint is non-binding (the optimiser naturally finds higher fill
rates) so those points do not lie on the true Pareto front; the trimmed
range $\geq 92\%$ ensures the constraint actually binds.


%% -----------------------------------------------------------------------
\section{Re-Evaluation and Search Space Summary}
\label{sec:opt:reeval}

Optuna internally minimises a noisy objective computed from a single
simulation run per trial.  Because the demand series is deterministic
(seeded) in these experiments, the noise comes only from the trial-to-trial
variation in dispatch decisions under different threshold values.  However,
for statistical rigour, the best parameter set identified by the optimiser
is \emph{re-evaluated} in a fresh simulation run before its cost and fill
rate are reported.

The re-evaluated metrics — not the optimiser's internal trial value —
are what appear in Chapter~\ref{chap:results}.  The best parameters are
also saved to a params.json file in the experiment output
directory, enabling reproducibility without re-running the optimiser.

%% -----------------------------------------------------------------------
\subsection{Search space dimensionality}
\label{sec:opt:dims}

Table~\ref{tab:opt-dims} summarises the search space for each optimisation
mode on the 1N3 ($N=2$), 5-SKU network.

\begin{table}[htbp]
  \centering
  \caption{Search space dimensionality by optimisation mode (1N3 with $N=2$ warehouses, 5 SKUs).}
  \label{tab:opt-dims}
  \begin{tabular}{@{}lccc@{}}
    \toprule
    Mode & Integer params & Continuous params & Total \\
    \midrule
    Inventory-only  & 25 & 0  & 25 \\
    Transport-only  & 0  & 5  & 5  \\
    Joint           & 25 & 5  & 30 \\
    \midrule
    E6 base\_stock (joint)     & 25 & 5  & 30 \\
    E6 $(s,S)$ (joint)        & 25+5$\alpha$ & 5 & 35 \\
    E6 periodic\_review (joint)& 25+5$R$      & 5 & 35 \\
    E6 echelon\_stock (joint)  & 25 & 5  & 30 \\
    \bottomrule
  \end{tabular}
\end{table}

The 30-dimensional joint space is well within the practical range of
TPE-based Bayesian optimisation (commonly cited as up to $\sim\!50$
dimensions).  Empirically, per-trial wall-clock time is flat from 5 to
30 dimensions (127--140\,ms/trial), confirming that TPE bookkeeping
overhead is negligible relative to the 365-day simulation evaluation;
100 trials require approximately 13 seconds per experiment.  Full
per-experiment runtimes and peak memory measurements are reported in
Appendix~\ref{appn:params}.

\paragraph{Per-policy parameter registry (E6).}
Experiment~E6 (Section~\ref{sec:results:e6}) compares four policy
families, each jointly optimising its inventory parameters
\emph{and} all five dispatch thresholds, so that every policy has
equal access to transport consolidation benefits.  The per-policy
registry in optimizer/policy\_search.py extends the base inventory
space: $(s,S)$ adds an $\alpha = s/S$ ratio per node (one shared across SKUs
at each node); periodic-review adds a review period $R$ per node;
echelon base-stock reuses a single level.  All variants then include the five continuous
dispatch thresholds $D_{\text{lane}}$, giving effective search
spaces of 30 dimensions (base\_stock, echelon\_stock) and 35
dimensions (ss, periodic\_review).
\paragraph{Bullwhip-aware joint optimisation (E8).}
Experiment~E8 (Section~\ref{sec:results:e8}) augments the joint
objective with a bullwhip term:
\vspace{-0.1em}
\begin{equation}
  \min \;\; C_{\text{total}}
            + \lambda \cdot \overline{\text{BW}}
            + \text{penalty}(\hat{\alpha})
  \label{eq:bullwhip-objective}
\end{equation}

where $\overline{\text{BW}}$ is the mean bullwhip ratio and $\lambda$
trades off cost against variance amplification.

The three optimisation modes --- inventory-only, transport-only, and
joint --- together with the NSGA-II Pareto study, the per-SKU
constraint mechanism, and the bullwhip-aware objective, form a complete
optimisation framework.  Experimental results obtained using this
framework are presented in Chapter~\ref{chap:results}.
